

\fancyhf{} % clear all header and footer fields
\fancyhead[RO,R]{\thepage} %RO=right odd, RE=right even
\renewcommand{\headrulewidth}{0pt}

\begin{center}
    \Large
    \textbf{Sokoban: History and AI}

    \vspace{0.4cm}
    \large
    Thesis Subtitle

    \vspace{0.4cm}
    \textbf{Author Name}

    \vspace{0.9cm}
    \textbf{Abstract}
\end{center}

This thesis investigates whether neuroevolution can produce heuristic functions for Weighted
A* search on Sokoban that are drastically more parameter-efficient than a conventional
gradient-trained baseline. Sokoban is a combinatorial planning puzzle that is both NP-hard and
PSPACE-complete, and solving it at scale requires an accurate cost-to-go heuristic to guide
search. The standard approach fixes a neural network architecture and trains it by gradient
descent; this thesis asks whether evolutionary search over network topology can find heuristics
that are competitive in accuracy while using orders of magnitude fewer parameters.

Three heuristic approaches are implemented and compared on the Boxoban benchmark, a
dataset of 9 million procedurally generated $10\times10$ Sokoban levels released by Google.
The CNN baseline is a fixed convolutional architecture with 2 million parameters trained by
gradient descent on cost-to-go labels derived from Breadth-First Search solution trajectories.
NEAT (NeuroEvolution of Augmenting Topologies) evolves both the topology and weights of a
feed-forward heuristic from scratch, starting from approximately 26 parameters and growing
complexity only where it improves prediction accuracy. CoDeepNEAT coevolves populations of
network modules and blueprint topologies, combining evolutionary architecture search with
per-genome gradient-descent weight training, producing assembled networks ranging from
roughly 32,000 to 170,000 parameters.

All three approaches are evaluated by mean absolute error on cost-to-go prediction and by
solve rate under Weighted A* search on held-out Boxoban levels. The results establish the
degree to which neuroevolution can close the accuracy gap relative to the gradient-trained
baseline while maintaining a substantial advantage in parameter count.